% ============================================================
% Ejercicio 3
% ============================================================
\section{Ejercicio 3}

\begin{tcolorbox}[enunciado]
Supongamos que, debido a una promoción agresiva de la aplicación, todos los pedidos disponibles en el sistema tienen exactamente el mismo peso (por ejemplo, todos pesan 500 gramos), aunque la ganancia de cada uno sigue siendo diferente dependiendo de la distancia de entrega.

Describe e implementa un algoritmo eficiente para resolver el problema planteado, para este caso particular (pesos idénticos). La implementación del algoritmo debe hacer uso del esquema de codificación implementado en el inciso 2.

El programa implementado debe recibir como entrada (desde consola) el nombre del archivo que contiene el ejemplar a probar, y como salida deberá imprimir:
\begin{itemize}[leftmargin=*]
    \item Respuesta al problema de decisión planteado (SÍ o NO se puede alcanzar la ganancia meta sin romper la mochila).
    \item Salida adicional: Si la respuesta es SÍ, imprimir el peso total.
\end{itemize}

Incluir al menos 2 archivos de prueba (ejemplos de ejecución) para ambos programas (incisos 2 y 3), así como ejemplos visuales de los ejemplares con los que prueben su programa.

Al menos deben agregar ejemplos de los siguientes casos
\begin{itemize}[leftmargin=*]
    \item ejemplar en el que la respuesta sea un sí
    \item ejemplar en el que la respuesta sea un no
\end{itemize}

* En caso de que no se pueda construir alguno de los ejemplares, justifica detalladamente por qué y propón algún otro caso de prueba.
\end{tcolorbox}

\subsection{Solución}
\subsubsection*{1. Análisis}
Dado que cada pedido pesa exactamente lo mismo (llamémosle $w_{const}$), el límite de la mochila $W$ ya no restringe combinaciones complejas de pesos, sino simplemente la \textbf{cantidad máxima de pedidos} que se pueden llevar.

La cantidad máxima de paquetes $k$ que caben en la mochila se calcula con una división entera:
$$k = \lfloor \frac{W}{w_{const}} \rfloor$$

Al saber que solo podemos llevar a lo mucho $k$ pedidos y que todos consumen exactamente el mismo espacio, la estrategia óptima es seleccionar los $k$ pedidos que ofrezcan la mayor ganancia individual.

\subsubsection*{Diseño del Algoritmo Eficiente}


\textbf{Lectura de Datos:} Se lee el archivo utilizando el esquema que diseñamos en el Ejercicio 2. Se asume que todos los pedidos tienen el mismo peso, por lo que tomamos ese valor como $w_{const}$.

\textbf{Cálculo de Límite de Artículos:} Se determina cuántos pedidos caben calculando $k = \lfloor W / w_{const} \rfloor$. Si resulta que $k$ es mayor que el número total de pedidos disponibles $n$, ajustamos $k = n$ (el repartidor tiene espacio para llevar todo el catálogo).

\textbf{Ordenamiento (Sorting):} Se ordena la lista de ganancias de los pedidos en orden descendente (del más lucrativo al menos lucrativo).

\textbf{Selección:} Se toman los primeros $k$ elementos de esta lista ordenada y se suman para obtener la ganancia máxima posible ($V_{max}$).

\textbf{Evaluación Final:}
\begin{itemize}[leftmargin=*]
    \item Si $V_{max} \ge V$, el algoritmo imprime ``SÍ'' y reporta el peso total transportado, el cual será exactamente la cantidad de artículos seleccionados multiplicada por $w_{const}$.
    \item Si $V_{max} < V$, el algoritmo imprime ``NO''.
\end{itemize}


\textbf{Ventaja de Complejidad}
En este caso donde los pesos son idénticos, la complejidad de nuestro algoritmo está dominada únicamente por el paso de ordenamiento de las ganancias, lo que nos da un tiempo de ejecución de $O(n \log n)$. Por lo que se puede considerar un algoritmo ``eficiente''.

\subsubsection*{2. Implementación en Python}
La implementación del algoritmo se encuentra en el archivo \texttt{mochila\_pesos\_identicos.py}


\subsubsection*{3. Descripción de los Ejemplos Visuales de los Ejemplares}
\textbf{Ejemplar 1: Caso Afirmativo (Respuesta ``SÍ'')}

Para este caso, diseñaremos un catálogo donde el repartidor tiene espacio limitado, pero al aplicar nuestro algoritmo, los paquetes más rentables logran superar la ganancia meta.

\begin{itemize}[leftmargin=*]
    \item \textbf{Total de pedidos ($n$):} 4
    \item \textbf{Capacidad de la mochila ($W$):} 1200 gramos
    \item \textbf{Ganancia meta ($V$):} 150 pesos
\end{itemize}

\textbf{Representación visual (sin codificar):}
\begin{center}
\renewcommand{\arraystretch}{1.2}
\begin{tabular}{ccc}
\toprule
\textbf{Pedido} & \textbf{Peso} & \textbf{Ganancia} \\
\midrule
1 & 500g & 30 pesos \\
2 & 500g & 100 pesos \\
3 & 500g & 50 pesos \\
4 & 500g & 80 pesos \\
\bottomrule
\end{tabular}
\end{center}

\textbf{Cadena que codifica al ejemplar (\texttt{ejemplar\_si.txt}):}
\begin{center}
\begin{terminal}
4\\
1200\\
150\\
500 100\\
500 80\\
500 50\\
500 30
\end{terminal}
\end{center}

\textbf{Resultado de la ejecución:}\\
Al ejecutar el algoritmo calculará que en la mochila solo caben 2 pedidos enteros (1200 divido entre 500). Tomará las ganancias mayores (100 y 80), sumando 180 pesos. Como 180 es mayor a la meta de 150 pesos, la salida en consola será:

\begin{center}
    \begin{terminal}
    SÍ\\
    1000
    \end{terminal}
\end{center}

\textbf{Ejemplar 2: Caso Negativo (Respuesta ``NO'')}

Aquí plantearemos un escenario restrictivo. La mochila es pequeña y, aunque el repartidor tome el pedido más caro, no será suficiente para alcanzar su meta.

\begin{itemize}[leftmargin=*]
    \item \textbf{Total de pedidos ($n$):} 3
    \item \textbf{Capacidad de la mochila ($W$):} 800 gramos
    \item \textbf{Ganancia meta ($V$):} 120 pesos
\end{itemize}

\textbf{Representación visual (sin codificar):}
\begin{center}
\renewcommand{\arraystretch}{1.2}
\begin{tabular}{ccc}
\toprule
\textbf{Pedido} & \textbf{Peso} & \textbf{Ganancia} \\
\midrule
1 & 500g & 40 pesos \\
2 & 500g & 90 pesos \\
3 & 500g & 70 pesos \\
\bottomrule
\end{tabular}
\end{center}

\textbf{Cadena que codifica al ejemplar (\texttt{ejemplar\_no.txt}):}
\begin{center}
    \begin{terminal}
    3\\
    800\\
    120\\
    500 90\\
    500 70\\
    500 40
    \end{terminal}
\end{center}

\textbf{Resultado de la ejecución:}\\
La mochila solo puede transportar 1 pedido (800 dividido entre 500). El mejor pedido disponible ofrece 90 pesos. Dado que 90 es menor que la meta de 120 pesos, es imposible cumplir el objetivo. La salida en consola será:

\begin{center}
    \begin{terminal}
    NO
    \end{terminal}
\end{center}


\subsubsection*{4. Análisis del Segundo Algoritmo Propuesto de Bajo Consumo}
Se sabe que los pesos para cada pedido $p$ son de igual cantidad; de esta manera $p_{w1}=p_{w2}=p_{w3}.. $, por lo tanto, lo único que posee un valor independiente $p_v$ son los valores para cada $p$. Teniendo esto en cuenta, se podría generar una suma de las cantidades de $p_v$ hasta que se alcance un $V$ de ganancias, siempre y cuando no sobrepase $W$. Esto se puede lograr gracias a un $max-heap$, ya que podemos ordenar los valores de los pedidos dependiendo de la cantidad de ganancia que se pueda ir generando. Para esto se utilizó la biblioteca \textit{heapq} de Python, que ayuda a poder utilizar los montículos mínimos. \cite{python_heapq}

Una vez implementados los montículos para las tuplas de pedidos, se convirtieron los valores de ganancias de cada pedido en valores negativos, para que el montículo mínimo tomara los valores negativos como los mejores valores; después se transformaba a un entero positivo con la función $abs()$ que convierte en valor absoluto.

La complejidad en tiempo está dada por $O(nlog(m))$, donde

La decodificación y validación de pesos del ejercicio anterior es: $O(n)$, al recorrer la lista de pedidos una vez para verificar que todos tengan el mismo peso $w$. Construcción del Max-Heap $(heapify):$ $O(n)$, tiempo lineal para organizar los $n$ elementos en la estructura del montículo. Selección de los mejores elementos: $O(m \log n)$, ya que el bucle realiza como máximo $m$ extracciones, y cada reordenamiento en un montículo de tamaño $n$ toma $O(\log n)$. De esta manera, $O(m\log n)$.

La complejidad en espacio del método es $O(n)$, donde $n$ es la cantidad total de pedidos. Lista de pedidos: Almacena las $n$ tuplas (peso, ganancia) procesadas por desbinarizar ($O(n)$). Estructura max-heap: Crea una nueva lista con $n$ tuplas $(-ganancia, peso)$ para la cola de prioridad ($O(n)$). Variables auxiliares: Variables como maxsum, maxweight, W y V ocupan un espacio constante e independiente del tamaño de la entrada ($O(1)$). Así $O(n)$.

El ejemplar que se toman en cuenta, para el si tiene los datos $ejemplar3$:

\begin{verbatim}
    Nombre del archivo a guardar (ej. ejemplar.txt): ejemplar3
¿Cuántos pedidos deseas ingresar? 6
Pedido 1:
Peso (en gramos): 300
Ganancia (en pesos): 40
Pedido 2:
Peso (en gramos): 300
Ganancia (en pesos): 24
Pedido 3:
Peso (en gramos): 300
Ganancia (en pesos): 17
Pedido 4:
Peso (en gramos): 300
Ganancia (en pesos): 26
Pedido 5:
Peso (en gramos): 300
Ganancia (en pesos): 8
Pedido 6:
Peso (en gramos): 300
Ganancia (en pesos): 10
Capacidad máxima de la mochila W (en gramos): 900
Ganancia mínima requerida V (en pesos): 75

Cadena binaria generada:

00000001001011000000000000101000000000010010110000000000000110000000000
1001011000000000000010001000000010010110000000000000
11010000000010010110000000000000010000000000100101100000000000000
101000000011100001000000000001001011
\end{verbatim}

El ejemplar que se toman en cuenta, para el si tiene los datos $ejemplar2$:

\begin{verbatim}
    Nombre del archivo a guardar (ej. ejemplar.txt): ejemplar2
¿Cuántos pedidos deseas ingresar? 4
Pedido 1:
Peso (en gramos): 300
Ganancia (en pesos): 40
Pedido 2:
Peso (en gramos): 200
Ganancia (en pesos): 30
Pedido 3:
Peso (en gramos): 350
Ganancia (en pesos): 17
Pedido 4:
Peso (en gramos): 299
Ganancia (en pesos): 19.5
[Error] Entrada inválida. Por favor, ingresa un número entero.
Ganancia (en pesos): 19
Capacidad máxima de la mochila W (en gramos): 800
Ganancia mínima requerida V (en pesos): 88

Cadena binaria generada:
00000001001011000000000000
101000000000001100100000000000000
111100000000101011110000000000001000
100000001001010110000000000010011
00000011001000000000000001011000

\end{verbatim}

